\scanpage{344}
% R344.01 | text | top centre | 473
\mspage{473}
% R344.02 | text | heading | VII Récapitulation sur les épinglages
$\overline{\text{\uline{VII}}}$\quad \uline{Recapitulation on pinnings.}
% R344.03 | text | upper, lines 2--4 | On désigne par G,SG,PG
\deleted{\unsure{1)}} We denote by $G,SG,PG$ forms of $\mathrm G\ell(2)_S$, $\mathrm S\ell(2)_S$, $\unsure{\mathrm{GP}(1)}_S$ \unsure{associated} to the same Azumaya Algebra $\mathcal M$, i.e. to a bundle of projective lines $\unsure{\mathbb P|\mathcal P}$, over $S$.
% R344.04 | text | subsection heading | Sections régulières de G
1) \uline{Regular sections of} $G$.
% R344.05 | text | upper middle, line 6 | Soit rho une section
Let $\rho$ be a section of $\unsure{\mathcal M}$ \inserted{(not necessarily invertible)}, put
% R344.06 | display | upper middle, line 7 | v=det rho,w=Tr rho
\[
(1)\qquad \unsure{v|\nu}=\det\rho,\quad w=\operatorname{Tr}\rho
\]
% R344.07 | text | upper middle, line 8 | de sorte qu'on aura
so that we have\deleted{ura} the Hamilton-Cayley equation
% R344.08 | display | upper middle, line 9 | rho^2-w rho+v=0
\[
(2)\qquad \rho^2-w\rho+\unsure{v|\nu}=0\qquad\text{i.e.}\quad\rho^2=w\rho-\unsure{v|\nu}
\]
% R344.09 | text | middle, line 10 | Proposition Conditions équivalentes
\uline{Proposition} Equivalent conditions:
% R344.10 | text | middle, line 11 | a) rho n'est un scalaire
a) $\rho$ is not a scalar on any fiber.
% R344.11 | text | middle, lines 12--13 | a') A_rho=O_S+O_S rho
\unsure{a')} $\mathcal A_\rho=\underline{\mathcal O}_S+\underline{\mathcal O}_S\rho$ is a submodule \unsure{locally|locally} a free direct summand of rank $2$ of $\mathcal M$.
% R344.12 | margin note | left, beside conditions a' and b | NB À cause de (2)
\margin{N.B. Because of (2), $\mathcal A_\rho$ is a sub-\uline{Algebra} of $\mathcal M$.}
% R344.13 | text | middle, lines 14--16 | b) Il existe loc. (Zar)
b) (If $\mathcal M=\underline{\operatorname{End}}_{\mathcal O_S}(F)$) There exists locally (Zariski) a section $\unsure{e_0|E_0}$ of $F$ such that $(e_0,\deleted{\unsure{\rho.}}e_1\mathrel{\overset{\mathrm{dfn}}{=}}\rho(e_0))$ form a basis of $F$
% R344.14 | text | middle-lower, lines 17--18 | b') Il existe loc. Zar
b') (If $\mathcal M=\underline{\operatorname{End}}_{\mathcal O_S}(F)$) There exists locally Zariski a ``line subbundle'' $D_0\subset F$, such that
% R344.15 | display | middle-lower, line 19 | F=D_0+D_1
\[
F=D_0+D_1,\quad\text{where }D_1=\rho(D_0)
\]
% R344.16 | text | lower, lines 20--22 | b'') Le O_S[T]-module
\unsure{b'')} (If $\mathcal M=\underline{\operatorname{End}}_{\mathcal O_S}(F)$.) The $\mathcal O_S[T]$-module $F_\rho$ ($\unsure{:}$ $T$ \unsure{is|being} an indeterminate, which operates on $F$ by $\rho$) is \deleted{\unsure{isomorphic} to $\underline{\mathcal O}_S$} \inserted{\deleted{\unsure{locally}} cyclic} locally (Zariski).
% R344.17 | text | lower, lines 23--24 | c) Le O_S[T]-module
c) (If $\mathcal M=\underline{\operatorname{End}}_{\mathcal O_S}(F)$) The $\underline{\mathcal O}_S[T]$-module $F_\rho$ \unsure{is|$\to$} isomorphic to
% R344.18 | display | lower, line 24 | O_S[T]/(T^2-wT+v)
\[
\underline{\mathcal O}_S[T]/(T^2-wT+\unsure{v|\nu})
\]
% R344.19 | text | bottom, lines 25--26 | c') il existe localt Zar une base
c') \inserted{(case $\mathcal M=\underline{\operatorname{End}}(F)$)} $\exists$ \unsure{locally|locally} Zariski a basis $e_0,e_1$ of $F$, such that the matrix of $\rho$ \unsure{with respect|p.v.} to this basis
% R344.20 | text | bottom left | 343
\mspage{343}

\scanpage{345}

is
\[
(3)\qquad \rho_{e_0,e_1}=\begin{pmatrix}0&-\unsure{v|\nu}\\1&w\end{pmatrix}
\]
\margin{(i.e. $\rho e_0=e_1$ whence \deleted{$\rho e_1=$}
$\rho e_1=\unsure{-v\rho e_0-w e_0}$}

c'') \inserted{(Case $\mathcal M=\underline{\operatorname{End}}_{\underline{\mathcal O}_S}(F)$)}
$\exists$ locally Zariski a basis $(e_0,e_1)$ of $F$ such that the matrix of $\rho$ with respect to it is
\[
(4)\qquad \rho_{e_0,e_1}=\begin{pmatrix}0&1\\-v&w\end{pmatrix}
\]

d) The commutant of $\rho$ in $\mathcal M$ is reduced to
$\mathcal A_\rho=\underline{\mathcal O}_S+\underline{\mathcal O}_S\rho$, and it remains so after every base change $S'\to S$.

d') \unsure{Item|Item} in each fiber.

\uline{Proof}\quad $a\Longleftrightarrow a'$ trivial (linear algebra of the Modules $\mathcal M$ and of the submodules $\underline{\mathcal O}_S1$ locally direct summands)

$b\Longleftrightarrow b'\Longleftrightarrow b''$ trivial, $b)\Longrightarrow a)$ as well.

$a\Longrightarrow b$ it suffices to choose $e_0$ \deleted{in}\inserted{on} a \deleted{\unsure{field}}\inserted{fiber} so as not to be an eigenvector under $\rho$ (possible over a field, for a nonscalar operator) then to lift in a neighborhood.

\inserted{$b\Longleftrightarrow$}$c)\Longleftrightarrow c')$ trivial ; \deleted{$c')\Longleftrightarrow c'')$ \unsure{for the}}\inserted{the bases for which}
$\rho_{e_0,e_1}$ is of the form $\begin{pmatrix}0&-v\\1&w\end{pmatrix}$ \unsure{are} precisely the bases of the form $(e_0,e_1=\rho e_0)$.

$c')\Longleftrightarrow c'')$ means that if $\rho$ satisfies the conditions under consideration, ${}^t\rho$ does too. (\deleted{\unsure{therefore}} and conversely it follows from this that the condition for ${}^t\rho$ implies it for $\rho={}^t({}^t\rho)$\deleted{\unsure{$\}$}}, i.e. $c'')\longrightarrow c')$) --- but this is immediate in the form a).

Thus the conditions a) to c'') are equivalent, moreover $d)\Longrightarrow d')\Longrightarrow a)$ is trivial, it remains to prove that $b$ (for example) $\Rightarrow d)$.
\deleted{If $u$ commutes}\inserted{Since $b$ \unsure{stable}} under base change, it suffices to see that

\mspage{344}

\scanpage{346}
\mspage{474}

\deleted{\uline{Corollary 2}}

\uline{N.B.}\quad \unsure{In} fact, $G=\underline{\mathcal M}^{*}$ operates on $\underline{\mathcal M}$ (by inner automorphisms) or on $G$, or $SG$, via $PG=\underline{\mathcal M}^{*}/\mathbb G_m$, and $PG$ is an fppf quotient of $SG$ (and \unsure{even|=} for the topology (étale), if 2 invertible).

So we see that \inserted{in} the preceding conjugation criterion, in the formula
\[
(6)\qquad \rho'=\operatorname{int}(\alpha)\rho=\alpha\rho\alpha^{-1},
\]
\unsure{we may assume} $\alpha\in\Gamma\unsure{\mathrm S\ell(2)|SG}$.

\deleted{\uline{Proposition} Let $\unsure{\varepsilon}$ be}\inserted{\uline{Corollary 2} Let} $\rho,\unsure{\zeta|\xi}$ be two sections of $G$, $\rho$ being regular. For us to be able to find locally Zariski an $\alpha\in\Gamma G$ such that
\[
(7)\qquad [\alpha,\rho]=\zeta\qquad\text{i.e.}\quad \alpha\rho\alpha^{-1}=\zeta\rho
\]
it is necessary and sufficient that we have

a) $\zeta\rho$ is regular, i.e. $\zeta$ \deleted{$\unsure{\neq}$} non-homothetic to $\rho^{-1}$ on any fiber

b) $\det\zeta=1$, $\operatorname{Tr}\zeta\rho=\operatorname{Tr}\rho$

(We can then take $\alpha$ in $SG$\ldots)

\unsure{This is|These are} only the conditions of corollary 1 \deleted{\unsure{of}} for conjugation of $\zeta\rho$ to $\rho$, the condition $\det\zeta\rho=\det\rho$ also being written as $\det\zeta\det\rho=\det\rho$ i.e. $\det\zeta=1$.

\uline{Corollary 3} The indeterminacy of $\alpha$ in the solution of (7) is of the form $\alpha\mapsto\alpha\uline{t}$, with $\uline{t}$ a section of $\underline{\mathcal A}_\rho^*$, i.e. $t=c\rho+d$, with $c,d\in\Gamma\underline{\mathcal O}_S$, $\det(c\rho+d)$ invertible.
\inserted{i.e. $\boxed{\text{\textcircled{8}}\quad c^2\unsure{v|\nu}+wcd\deleted{\text{invertible}}+d^2\quad\text{invertible}}$}
(where, $v=\det\rho$, $w=\operatorname{Tr}\rho$, cf. (1)).

\mspage{345}

\scanpage{347}

if $u$ is a section of $\mathcal M$ \deleted{such that} which commutes with $\rho$, then $u=a+b\rho$ $(a,b\in\Gamma\underline{\mathcal O}_S)$. By a) $a$ and $b$ are unique, and this shows that the question is fppf local. \unsure{We may assume} therefore $\mathcal M=\underline{\operatorname{End}}(F)$, $F$ equipped with a basis $(e_0,e_1=\rho(e_0))$. \deleted{\unsure{and}} We will have, if
\[
\deleted{ue_1=u\rho e_0=\rho=}
\]
\[
u(e_0)=ae_0+be_1\inserted{{}=u'(e_0)}\qquad
\deleted{\text{put }u'=a+b\rho}\inserted{\text{where}}\quad u'=a+b\rho,
\]
\[
\deleted{u(e_1)=u\rho(e_0)=\rho\underline{ue_0}=\rho(ae_0+be_1)=a\rho(e_0)+b\rho(e_1)}
\]
and if $u$ \unsure{commutes} \deleted{$=ae_1+$} with $\rho$, we will have
\[
u(e_1)=u\rho(e_0)=\rho u(e_0)\deleted{{}=\rho(ae_0+}
\]
likewise
\[
u'(e_1)=u'\rho(e_0)=\rho u'(e_0)\qquad\text{for }u'\text{ commutes with }\rho
\]
therefore
\[
u(e_0)\unsure{=|\neq}\unsure{u(e'_0)|u'(e_0)},\quad u(e_1)=u'(e_1)\qquad\text{therefore }u=u',
\]
q.e.d.

\uline{Definition :} A section \inserted{of $\mathcal M$} which satisfies the equivalent conditions of the proposition is called \uline{regular}. \unsure{Item|Item} for sections of $G$, $SG$, interpreted as sections of $\mathcal M$ --- \deleted{they are} such a section is regular iff at no point is it a scalar, i.e. is not in the center. This \unsure{remains|is} stable under multiplication by an invertible scalar --- whence by passage to the quotient we deduce the notion of a regular section of $PG$ --- these are the sections which are not $=1$ on any fiber.

\uline{Corollary 1} Let $\rho,\rho'$ be two \deleted{regular} sections of $\mathcal M$, \inserted{$\rho$ being regular.} For them to be conjugate (locally Zariski) via $\mathcal M^*$, it is necessary and sufficient that \deleted{$\rho$} a) $\rho'$ be regular (Condition \inserted{of regularity} \deleted{\illegible{1}}) and b) we have
\[
(5)\qquad \det\rho=\det\rho',\quad \operatorname{Tr}\rho=\operatorname{Tr}\rho'
\]
(``Numerical condition''

\mspage{346}

\scanpage{348}
\mspage{475}

This follows from the equivalence $a)\Longleftrightarrow d)$ of the proposition.

\uline{Proposition} Let
\[
(9)\qquad \varepsilon=1+N
\]
be a section of $\mathcal M$. Equivalent conditions

\[
\text{a)}\quad (10)\quad
\begin{gathered}
\det\varepsilon=\det1\quad(=1)\\
\operatorname{Tr}\varepsilon=\operatorname{Tr}1\quad(=2)
\end{gathered}
\qquad \deleted{\varepsilon\text{ regular}}\quad\text{and }\varepsilon\text{ regular}
\]
\margin{\uline{N.B.}\quad $\det(1+N)=1+\operatorname{Tr}N+\det N$\quad $\operatorname{Tr}(1+N)=2+\operatorname{Tr}N$.}
\[
\text{b)}\quad (11)\quad \det N=\operatorname{Tr}N=0\qquad\text{and }N\text{ regular}
\]

c) (If $\mathcal M=\underline{\operatorname{End}}_{\underline{\mathcal O}_S}(F)$) $\exists$ line bundle $D\subset F$ such that we have
\[
N(F)=D,\quad N(D)=0
\]
(Then $D$ is uniquely determined by $\varepsilon$ i.e. by $N$, \deleted{and conversely} \inserted{$D$ determines $N$ modulo homothety})

\inserted{c') (If $\mathcal M=\underline{\operatorname{End}}_{\underline{\mathcal O}_S}(F)$). $\exists$ locally \unsure{basis|Zariski} $e_0,e_1$ such that $\varepsilon=\begin{pmatrix}1&1\\0&1\end{pmatrix}$ i.e. $N=\begin{pmatrix}0&1\\0&0\end{pmatrix}$}

d) $\exists$ Borel $B\subset G$, such that $\varepsilon$ is a section everywhere $\ne1$ of $B_u$.

d') \deleted{\unsure{If}}$\exists$ Borel $SB\subset SG$, such that $\varepsilon$ is a section everywhere $\ne1$ of $(SB)_u$.

\deleted{\unsure{Proof}} Moreover, $D,B,SB$ are well determined by $\varepsilon$ i.e. by $N$, and determine $N$ up to homothety.

\uline{Definition :} we \deleted{say} \unsure{that} $\varepsilon$ is (strictly) unipotent iff we have $\det\varepsilon=1$, $\operatorname{Tr}\varepsilon=2$, that $N$ is \uline{strictly nilpotent} \unsure{iff|if} $\det N=\operatorname{tr}N=0$ (whence $N^2=0$ by Hamilton-Cayley). We say that $\varepsilon$ is regular unipotent resp. regular nilpotent (dropping the qualifier ``strictly'') if in addition it is regular.

\mspage{347}